\chapter[Additional Technical Details]{Additional Technical\\Details}\label{app-phase}
\section{Uniform phase convention}
The overall sign of an electronic wave function is arbitrary. Multiplying one quasi-diabatic state by \(-1\) leaves its energy unchanged, but reverses the signs of all electronic couplings and transition properties that connect that state to other states. A reported coupling sign is therefore meaningful only after a phase convention has been specified.

All diabatization methods in \progname{} use a common fixed-geometry convention~\cite{Diabat2026}. After the adiabatic-to-diabatic transformation has been obtained, each quasi-diabatic state is compared with a phase-reference state. The program evaluates a transition dipole after symmetric orthogonalization of the two states. Only the alpha-spin contribution is used for the phase assignment. This avoids exact alpha-beta cancellation in cases such as an \(M_S=0\) triplet. The Cartesian component with the largest absolute value is identified. If that component is negative, the sign of the quasi-diabatic state is reversed.

When the diabatization method already uses a suitable reference state, that state is also used as the phase reference. If no such state is available, if the required transition dipole cannot be evaluated, or if the norm of the corresponding alpha-spin transition dipole is below \(10^{-8}\) atomic units, the lowest-energy adiabatic state in the selected diabatization space is used instead. This convention makes repeated calculations at the same nuclear geometry and Cartesian frame deterministic and gives mutually consistent signs among the couplings produced by different supported diabatization methods.

The fixed-geometry phase convention is not a phase-tracking algorithm along nuclear motion. Changing the nuclear geometry changes the wave functions that are being compared. Rotating the Cartesian frame can also change the component used for the sign choice. Use Phasefix in Section~\ref{sec-phasefix} when target wave functions at one geometry need to be aligned with reference wave functions from another geometry.

\section{Reference states}
A reference state defines the baseline against which a transition or difference density is interpreted. The exact role depends on the calculation. In single-excitation FPHD, particle and hole information is obtained from transition densities between the target states and a reference. For the supported single-excitation wave-function representation, the underlying reference ground state is used automatically when no explicit reference is supplied. A loaded state outside the target space can be selected explicitly when a different supported reference is needed.

Multiple-excitation FPHD uses difference density matrices. The reference may therefore come from a different electronic-structure method and may have a different number of electrons. A common example is a cationic target manifold referenced to the neutral ground state so that charge localization and excitation character can be interpreted relative to the neutral system. The reference and target states must still share the same molecular geometry and compatible AO basis. This condition is what allows their density matrices to be represented and compared consistently.

FED also relies on excitation densities relative to a reference state. The standalone implementation permits two adiabatic target states. As with other difference-density analyses, a physically meaningful reference and a compatible one-electron representation are required. The detailed keyword behavior for reference selection is given with the corresponding job controls rather than being inferred from this general discussion.

\section{State classification and partial diagonalization}
Property-based diabatization often determines classes of states more robustly than it determines every individual state inside a class. Two intermediate quasi-diabatic states can have nearly the same descriptor while differing mainly by an arbitrary unitary rotation within their common subspace. In this situation, insisting that the property descriptor alone fix their individual orientation can make the result numerically unstable and chemically uninformative.

\progname{} therefore classifies the intermediate states according to the descriptors of the selected method. It then partially diagonalizes the Hamiltonian within each class when the corresponding calculation requires this treatment. The operation preserves the separation between classes while restoring an adiabatic representation inside a class that the descriptor cannot uniquely resolve. FPHD uses fragment particle and hole populations for classification. Multistate GMH uses the projected dipole descriptor. Multistate FCD uses the charge-difference descriptor. Multistate FED--FCD uses its combined charge and excitation descriptors.

The \keyref{key-common-partial_diag}{partial\_diag} keyword controls this final step where supported. The default behavior is designed for the implemented classification procedure. Disable or force it only when there is a specific reason to inspect the intermediate representation or to reproduce a deliberately modified protocol.

\section{Combining states from separate external calculations}
A \texttt{statepack} may load states from more than one external calculation. Whether those states can be combined meaningfully depends on the operation to be performed. Reading several files is not by itself a guarantee that their wave functions are compatible. Tasks involving transition or difference density matrices require consistent atom definitions and AO basis information. Transition-density calculations impose additional wave-function compatibility requirements.

For practical use, preserve the same molecular geometry, atom order, basis definition, and compatible orbital representation whenever separate calculations are intended to enter one density-matrix operation. The permitted exception in multiple-excitation FPHD concerns the electronic method and electron number of the reference state. It does not remove the shared geometry and AO-basis requirement. If a cross-file calculation fails, the program diagnostics should be considered together with the requirements of the specific method described in Chapters~\ref{chap-interfaces} and~\ref{chap-diabat}.
